\chapter{Séries de Fourier}\label{ch:b2:fourier}

Tout signal périodique peut-il être reconstruit à partir de purs
sinus et cosinus~? Le « oui » audacieux de Fourier a engendré un
siècle d'analyse. Ce chapitre démontre les deux piliers à portée de
la MP*~: le \emph{théorème de Dirichlet} (reconstruction ponctuelle
des fonctions $C^1$ par morceaux, via le \omterm{lem:b2:fourier:kernel}{noyau de Dirichlet}) et
l'\emph{identité de Parseval} (l'énergie d'un signal est la somme des
énergies de ses harmoniques), et récolte les séries numériques
classiques --- la somme de Bâle $\sum 1/n^2 = \pi^2/6$ en tête.

Dans tout le chapitre, les fonctions sont $2\pi$-périodiques,
\omterm{def:b2:metric:continuity}{continues} par morceaux, à valeurs complexes~; $\mathcal{C}$ désigne
les fonctions \omterm{def:b2:metric:continuity}{continues}.

\section{Coefficients de Fourier}

\begin{definition}\label{def:b2:fourier:coefficients}
Les \emph{coefficients de Fourier}\index{coefficients de Fourier} de
$f$ sont
\[
c_n(f) = \frac{1}{2\pi}\int_{-\pi}^{\pi} f(t)\,\eu^{-\iu n t}\,\dd
t \qquad (n \in \Z),
\]
et les coefficients sous forme réelle $a_n = c_n + c_{-n}$, $b_n =
\iu(c_n - c_{-n})$, de sorte que les \emph{sommes partielles de
Fourier}\index{sommes partielles de Fourier} sont
\[
S_N(f)(t) = \sum_{n=-N}^{N} c_n(f)\,\eu^{\iu nt}
= \frac{a_0}{2} + \sum_{n=1}^{N}\bigl(a_n\cos nt + b_n \sin
nt\bigr).
\]
Sur $\mathcal{C}$, on définit le \omterm{def:b2:hermitian:def}{produit scalaire hermitien} $\langle
f, g\rangle = \frac{1}{2\pi}\int_{-\pi}^{\pi}\conj f\,g$~: les
exponentielles $e_n(t) = \eu^{\iu nt}$ sont \emph{orthonormées}
($\langle e_m, e_n\rangle = \delta_{mn}$, calcul direct), et
$c_n(f) = \langle e_n, f\rangle$~: l'analyse de Fourier est de la
géométrie hermitienne (\cref{ch:b2:hermitian}) en dimension infinie.
\end{definition}

\begin{proposition}[Inégalité de Bessel]\label{prop:b2:fourier:bessel}
$S_N(f)$ est la projection orthogonale de $f$ sur l'espace
$\mathcal{T}_N$ des polynômes trigonométriques de degré $\leq N$,
et\index{inégalité de Bessel}
\[
\sum_{n=-N}^{N} \abs{c_n(f)}^2 \leq \norm f_2^2 =
\frac{1}{2\pi}\int_{-\pi}^{\pi} \abs f^2 :
\]
la série $\sum \abs{c_n}^2$ \omterm{def:b2:integration:improper}{converge}, et $c_n(f) \to 0$ quand
$\abs n \to \infty$ (Riemann--Lebesgue pour les coefficients).
\end{proposition}

\begin{proof}
$f - S_N(f)$ est orthogonal à chaque $e_k$, $\abs k \leq N$ ($\langle
e_k, f - S_N f\rangle = c_k - c_k = 0$)~: $S_Nf$ est la projection
orthogonale sur $\mathcal{T}_N = \operatorname{Vect}(e_{-N}, \dots,
e_N)$ (le théorème de projection du volume de Licence~1, mot pour mot
dans le cadre \omterm{def:b2:hermitian:adjoint}{hermitien}). Pythagore~: $\norm f_2^2 = \norm{S_Nf}_2^2 +
\norm{f - S_Nf}_2^2 \geq \norm{S_Nf}_2^2 = \sum_{\abs n \leq N}
\abs{c_n}^2$~; on fait $N \to \infty$.
\end{proof}

\begin{example}[Meilleure approximation, mesurée]\label{ex:b2:fourier:bestapprox}
Dans quelle mesure les polynômes trigonométriques de bas degré
approchent-ils la dent de scie $f(t) = t$ (sur $\intoo{-\pi}{\pi}$)
en moyenne quadratique~? D'après la~\cref{prop:b2:fourier:bessel}, la
meilleure approximation de degré $N$ \emph{est} $S_N(f)$, avec pour
erreur quadratique
\[
\norm{f - S_Nf}_2^2 = \norm f_2^2 -
\sum_{\abs n\leq N}\abs{c_n}^2 .
\]
Ici $\norm f_2^2 = \frac{1}{2\pi}\int_{-\pi}^\pi t^2\dd t =
\frac{\pi^2}{3}$, et à partir de $b_n = \frac{2(-1)^{n+1}}{n}$
(\cref{ex:b2:fourier:basel})~: $\abs{c_n}^2 + \abs{c_{-n}}^2 =
\frac{b_n^2}{2} = \frac{2}{n^2}$. Donc
\[
\norm{f - S_Nf}_2^2
= \frac{\pi^2}{3} - \sum_{n=1}^{N}\frac{2}{n^2}
\qquad\text{~: numériquement } 1.29,\ 0.79,\ 0.57,\ 0.44
\]
pour $N = 1, 2, 3, 4$ --- décroissante, mais lentement~: le reste
$\sum_{n>N}\frac2{n^2} \sim \frac2N$ est gouverné par la lente
décroissance en $\frac1n$ des coefficients, elle-même la signature
du saut (\cref{exo:b2:fourier:6} lu à l'envers). Éclairage final~:
Parseval transforme la qualité d'approximation en un reste de série
numérique --- et prédit, avant tout dessin, que les sauts font
converger les séries de Fourier à contrecœur.
\end{example}

\begin{method}[Calculer efficacement les coefficients de Fourier]\label{met:b2:fourier:coefficients}
Avant d'intégrer quoi que ce soit~:
\begin{enumerate}
  \item \emph{Parité~:} une fonction $f$ paire a $b_n = 0$, une
        fonction $f$ impaire a $a_n = 0$, et les intégrales
        survivantes se ramènent à $\frac2\pi \int_0^\pi$ --- moitié
        moins de travail, deux fois plus de fiabilité.
  \item \emph{Les polynômes trigonométriques sont déjà faits~:}
        linéariser les produits ($\cos^3$, $\sin^2\cos$, \dots) et
        lire directement les coefficients
        (\cref{exo:b2:fourier:9})~; l'orthonormalité rend toute
        intégration supplémentaire redondante.
  \item \emph{Exponentielles complexes pour les exponentielles~:}
        pour des facteurs $\eu^{at}$ ou des oscillations amorties,
        calculer $c_n$ directement --- une intégrale de $\eu^{(a -
        \iu n)t}$ vaut mieux que deux intégrations par parties
        (\cref{exo:b2:fourier:10}).
  \item \emph{Dériver un développement connu~:} si $f'$ a des
        coefficients connus et $f$ est \omterm{def:b2:metric:continuity}{continue}, $c_n(f) =
        \frac{c_n(f')}{\iu n}$ ($n \neq 0$) récupère tout sauf
        $c_0$, qui est la moyenne --- souvent la voie la plus
        rapide, et légitime exactement sous les hypothèses du
        \cref{thm:b2:fourier:parseval}~(1).
\end{enumerate}
\end{method}

\section{Théorème de Dirichlet}

\begin{lemma}[Noyau de Dirichlet]\label{lem:b2:fourier:kernel}
$S_N(f)(x) = \frac{1}{2\pi}\int_{-\pi}^{\pi} f(x + u)\,D_N(u)\,\dd
u$, où\index{noyau de Dirichlet}
\[
D_N(u) = \sum_{n=-N}^{N} \eu^{\iu nu}
= \frac{\sin\bigl((N + \frac12)u\bigr)}{\sin\frac u2}
\quad (u \notin 2\pi\Z),
\qquad
\frac{1}{2\pi}\int_{-\pi}^{\pi} D_N = 1 .
\]
\end{lemma}

\begin{proof}
On insère la définition de $c_n$ dans $S_N$ et on échange somme et
intégrale (légitime~: la somme est finie)~:
\[
S_N(f)(x)
= \sum_{n=-N}^{N}\Bigl(\frac{1}{2\pi}\int_{-\pi}^{\pi}
f(t)\,\eu^{-\iu nt}\dd t\Bigr)\eu^{\iu nx}
= \frac{1}{2\pi}\int_{-\pi}^{\pi} f(t)
\sum_{n=-N}^{N}\eu^{\iu n(x-t)}\,\dd t ;
\]
on substitue $u = t - x$ et on ramène le segment d'intégration à
$\intcc{-\pi}{\pi}$ à l'aide de la $2\pi$-périodicité de l'intégrande~;
la symétrie de la plage d'indices donne $\sum_n\eu^{-\iu nu} =
D_N(u)$. La forme fermée~: somme géométrique de raison $\eu^{\iu u}$,
\[
D_N(u) = \eu^{-\iu Nu}\,\frac{\eu^{\iu(2N+1)u} - 1}{\eu^{\iu u} -
1}
= \frac{\eu^{\iu(N + \frac12)u} - \eu^{-\iu(N+\frac12)u}}
{\eu^{\iu u/2} - \eu^{-\iu u/2}} ,
\]
qui est le quotient de sinus. Sa moyenne vaut $1$~: seul $n = 0$
contribue.
\end{proof}

\begin{theorem}[Lemme de Riemann--Lebesgue]\label{thm:b2:fourier:riemannlebesgue}
Pour $g$ \omterm{def:b2:metric:continuity}{continue} par morceaux sur un segment,
$\int_a^b g(t)\sin(\lambda t + \varphi)\,\dd t \to 0$ quand $\lambda
\to +\infty$.\index{lemme de Riemann--Lebesgue}
\end{theorem}

\begin{proof}
On approche $g$ \omterm{def:b2:funcseq:def}{uniformément} par des fonctions en escalier
(\cref{thm:b2:funcseq:weierstrass} n'est pas nécessaire --- l'approche
élémentaire par fonctions en escalier des fonctions \omterm{def:b2:metric:continuity}{continues} par
morceaux suffit) et on intègre chaque marche explicitement~: chaque
morceau contribue pour $O\bigl(\frac1\lambda\bigr)$, et l'erreur
d'approximation contribue pour $\varepsilon(b - a)$. Cet argument a
été mené en détail dans le tout dernier exercice du chapitre
d'intégration du volume de Licence~1~; pour des morceaux $C^1$ on peut
aussi intégrer par parties et majorer par $\frac C\lambda$.
\end{proof}

\begin{example}[À quelle vitesse les coefficients meurent-ils~?]\label{ex:b2:fourier:decaytable}
Riemann--Lebesgue dit que les coefficients tendent vers $0$~; leur
\emph{vitesse} est un instrument de mesure de la régularité. Trois
spécimens de ce chapitre et de ses exercices~:
\[
\text{signal carré~: } b_n = \frac{4}{\pi n}\ (n\ \text{impair}),
\qquad
\abs t : a_n = \frac{-4}{\pi n^2}\ (n\ \text{impair}),
\qquad
\abs{\sin t} : a_{2k} = \frac{-4}{\pi(4k^2-1)} .
\]
Un saut de $f$ (signal carré, dent de scie) laisse des coefficients
d'\omterm{def:b2:structures:generated}{ordre} $\frac1n$~: pas de \omterm{def:b2:funcseq:series}{convergence normale}, dépassement de Gibbs
aux sauts. La continuité avec un point anguleux --- un saut de $f'$
seulement --- améliore l'\omterm{def:b2:structures:generated}{ordre} à $\frac{1}{n^2}$~: \omterm{def:b2:funcseq:series}{convergence
normale}, reconstruction uniforme. En général $k$ dérivées achètent
$c_n = O(n^{-k})$ (\cref{exo:b2:fourier:6}), et réciproquement un
\omterm{def:b2:reduction:eigen}{spectre} décroissant plus vite que toute puissance force $f$ à être
$C^\infty$ (on dérive terme à terme, maintenant légitimement).
Éclairage final~: la régularité du signal et la décroissance du
\omterm{def:b2:reduction:eigen}{spectre} sont la même information --- un ingénieur lit l'une sur la
pente de l'autre sans jamais tracer la fonction.
\end{example}

\begin{theorem}[Dirichlet]\label{thm:b2:fourier:dirichlet}
Soit $f$ une fonction $2\pi$-périodique et $C^1$ par morceaux. Alors
pour tout $x$,
\[
S_N(f)(x) \xrightarrow[N \to \infty]{}
\frac{f(x^+) + f(x^-)}{2}
\]
(la moyenne des limites unilatérales) --- en particulier $S_N(f)(x)
\to f(x)$ en tout point de continuité.\index{théorème de Dirichlet}
\end{theorem}


\begin{proof}
Par le lemme du noyau et sa moyenne unité, en scindant l'intégrale en
les deux moitiés $u > 0$ et $u < 0$ (chacune de moyenne
$\frac12$)~:
\begin{align*}
S_N(f)(x) - \frac{f(x^+) + f(x^-)}{2}
&= \frac{1}{2\pi}\int_{0}^{\pi} \bigl(f(x+u) -
f(x^+)\bigr)D_N(u)\,\dd u\\
&\quad+ \frac{1}{2\pi}\int_{-\pi}^{0}\bigl(f(x+u) -
f(x^-)\bigr)D_N(u)\,\dd u .
\end{align*}
Traitons la première (la seconde est \omterm{def:b2:quadratic:adjoint}{symétrique}). Écrivons
\[
\bigl(f(x + u) - f(x^+)\bigr)\,D_N(u)
= \underbrace{\frac{f(x+u) - f(x^+)}{\sin\frac u2}}_{g(u)}\,
\sin\Bigl(\Bigl(N + \frac12\Bigr)u\Bigr) .
\]
La fonction $g$ est \omterm{def:b2:metric:continuity}{continue} par morceaux sur $\intoc{0}{\pi}$ et a
une \emph{limite finie en $0^+$}~: en écrivant
\[
g(u) = \frac{f(x+u) - f(x^+)}{u}\cdot\frac{u}{\sin\frac u2} ,
\]
le premier facteur tend vers $f'(x^+)$ (dérivabilité à droite,
conséquence du caractère $C^1$ par morceaux) et le second vers $2$
(la limite classique $\frac{\sin v}{v} \to 1$ en $v = \frac
u2$)~: $g(0^+) = 2f'(x^+)$ existe. Donc $g$ se prolonge de manière
\omterm{def:b2:metric:continuity}{continue} par morceaux à $\intcc{0}{\pi}$, et Riemann--Lebesgue
(\cref{thm:b2:fourier:riemannlebesgue}) envoie l'intégrale vers
$0$. C'est là tout l'enjeu de l'hypothèse~: sans dérivées
unilatérales, le facteur $\frac{1}{\sin(u/2)}$ explose en $0$ plus
vite que Riemann--Lebesgue ne peut compenser, et la \omterm{def:b2:funcseq:def}{convergence
simple} peut réellement échouer pour une fonction $f$ seulement
\omterm{def:b2:metric:continuity}{continue} --- la brèche que le théorème de Fejér (devoir) comble par
moyennisation.
\end{proof}

\begin{example}[Dirichlet en un saut]\label{ex:b2:fourier:jumpvalue}
Pour la dent de scie $f(t) = t$ sur $\intoo{-\pi}{\pi}$
(\cref{ex:b2:fourier:basel} ci-dessous), le prolongement périodique
saute en $t = \pi$ de $f(\pi^-) = \pi$ à $f(\pi^+) = -\pi$.
Dirichlet promet la valeur $\frac{\pi + (-\pi)}{2} = 0$ là, et en
effet chaque terme de $\sum \frac{2(-1)^{n+1}}{n}\sin nt$ s'annule
en $t = \pi$~: la série \omterm{def:b2:integration:improper}{converge} poliment vers le milieu, ignorant
les deux valeurs unilatérales. En déplaçant le point d'évaluation en
$t = \frac\pi2$ (un point de continuité), la même série devient le
$\frac\pi4$ de Leibniz. Une série, deux comportements --- exactement
les deux clauses du théorème.
\end{example}

\begin{theorem}[Convergence normale pour $C^1$~; Parseval]\label{thm:b2:fourier:parseval}
\begin{enumerate}
  \item Si $f$ est \omterm{def:b2:metric:continuity}{continue}, $2\pi$-périodique et $C^1$ par
        morceaux, alors $c_n(f') = \iu n\,c_n(f)$, la série de
        Fourier de $f$ \omterm{def:b2:integration:improper}{converge} \emph{\omterm{def:b2:funcseq:series}{normalement}} sur $\R$, et sa
        somme est $f$.
  \item (Parseval) Pour toute fonction $f$ $2\pi$-périodique
        \omterm{def:b2:metric:continuity}{continue} par morceaux~:\index{identité de Parseval}
        \[
        \frac{1}{2\pi}\int_{-\pi}^{\pi}\abs{f}^2
        = \sum_{n=-\infty}^{\infty} \abs{c_n(f)}^2
        = \frac{\abs{a_0}^2}{4} + \frac12\sum_{n\geq1}
        \bigl(\abs{a_n}^2 + \abs{b_n}^2\bigr).
        \]
        \emph{(Démontré ici pour $f$ \omterm{def:b2:metric:continuity}{continue} et $C^1$ par
        morceaux~; admis en général.)}
\end{enumerate}
\end{theorem}

\begin{proof}
(1) Intégration par parties sur chaque morceau $C^1$ (les termes de
bord se compensent par continuité et périodicité)~: $c_n(f') = \iu n
c_n(f)$. Puis, par Cauchy--Schwarz sur les deux familles de carré
\omterm{def:b2:series:summable}{sommable} (\cref{prop:b2:fourier:bessel} pour $f'$)~:
\[
\sum_{n \neq 0} \abs{c_n(f)} = \sum_{n\neq0}
\frac{\abs{c_n(f')}}{\abs n}
\leq \Bigl(\sum \abs{c_n(f')}^2\Bigr)^{1/2}
\Bigl(\sum_{n\neq0}\frac{1}{n^2}\Bigr)^{1/2} < \infty :
\]
\omterm{def:b2:funcseq:series}{convergence normale} de la série de Fourier. Sa somme est \omterm{def:b2:metric:continuity}{continue} et
coïncide avec $f$ en tout point d'après Dirichlet
(\cref{thm:b2:fourier:dirichlet}~: $f$ est \omterm{def:b2:metric:continuity}{continue})~: la série
\omterm{def:b2:integration:improper}{converge} vers $f$, \omterm{def:b2:funcseq:def}{uniformément}.

(2) Pour une telle $f$~: $S_N f \to f$ \omterm{def:b2:funcseq:def}{uniformément}, donc $\norm{f -
S_Nf}_2 \leq \norm{f - S_Nf}_\infty \to 0$, et Pythagore
($\norm f_2^2 = \sum_{\abs n \leq N}\abs{c_n}^2 + \norm{f -
S_Nf}_2^2$) passe à la limite. La forme réelle est de la
comptabilité avec $a_n, b_n$.
\end{proof}

\begin{example}[La majoration du reste $C^1$, rendue quantitative]\label{ex:b2:fourier:tailbound}
La démonstration du \cref{thm:b2:fourier:parseval}~(1) cache une
estimation utilisable. Pour $f$ \omterm{def:b2:metric:continuity}{continue} et $C^1$ par morceaux, le
même Cauchy--Schwarz appliqué au seul reste donne
\[
\sum_{\abs n > N}\abs{c_n(f)}
= \sum_{\abs n > N}\frac{\abs{c_n(f')}}{\abs n}
\leq \Bigl(\sum_{\abs n > N}\abs{c_n(f')}^2\Bigr)^{\!1/2}
\Bigl(\sum_{\abs n>N}\frac{1}{n^2}\Bigr)^{\!1/2}
\leq \norm{f'}_2\,\sqrt{\frac{2}{N}} ,
\]
en utilisant Bessel pour $f'$ et $\sum_{n>N}n^{-2} \leq \frac1N$.
Ainsi l'erreur uniforme des sommes partielles obéit à
\[
\norm{f - S_Nf}_\infty
\leq \sum_{\abs n>N}\abs{c_n(f)}
\leq \norm{f'}_2\,\sqrt{\frac2N} .
\]
Pour $f(t) = \abs t$~: $\norm{f'}_2 = 1$ (la dérivée est $\pm1$),
donc dix termes reconstruisent déjà $\abs t$ \omterm{def:b2:funcseq:def}{uniformément} à
$\sqrt{0.2} \approx 0.45$ près, et $N = 10^4$ à $0.015$ près.
Éclairage final~: une dérivée achète la vitesse uniforme en
$\frac{1}{\sqrt N}$~; en comparant avec le monde en $\frac1n$ des
coefficients du signal carré (aucune \omterm{def:b2:funcseq:def}{convergence uniforme}), le
dictionnaire de l'\cref{ex:b2:fourier:decaytable} se dote de
nombres.
\end{example}

\begin{example}[Bâle et compagnie]\label{ex:b2:fourier:basel}
Soit $f(t) = t$ sur $\intoo{-\pi}{\pi}$, prolongée de manière
$2\pi$-périodique (une dent de scie, $C^1$ par morceaux). En
calculant, $a_n = 0$ (imparité) et
\[
b_n = \frac{1}{\pi}\int_{-\pi}^{\pi} t\sin nt\,\dd t
= \frac{2(-1)^{n+1}}{n} .
\]
Dirichlet en $t = \frac\pi2$ redonne le $\frac\pi4 = 1 - \frac13 +
\frac15 - \dots$ de Leibniz~; Parseval donne
\[
\frac{1}{2\pi}\int_{-\pi}^{\pi} t^2\,\dd t = \frac{\pi^2}{3}
= \frac12\sum_{n\geq1}\frac{4}{n^2}
\quad\Longrightarrow\quad
\boxed{\;\sum_{n\geq1}\frac{1}{n^2} = \frac{\pi^2}{6}\;}
\]
--- la somme de Bâle d'Euler, en deux lignes. La fonction $f(t) =
t^2$ donne de même $\sum \frac1{n^4} = \frac{\pi^4}{90}$
(\cref{exo:b2:fourier:3}).
\end{example}

\begin{example}[Un développement complet avec vérification intégrée~: $\abs{\sin t}$]\label{ex:b2:fourier:abssin}
La fonction $f(t) = \abs{\sin t}$ est \omterm{def:b2:metric:continuity}{continue}, paire,
$\pi$-périodique (donc $2\pi$-périodique), $C^1$ par morceaux. La
parité tue les $b_n$~; $a_0 = \frac1\pi\int_{-\pi}^{\pi} \abs{\sin
t}\dd t = \frac4\pi$~; et pour $n \geq 1$, la transformation
produit-somme donne
\[
a_n = \frac2\pi\int_0^\pi \sin t\cos nt\,\dd t
= \frac{1}{\pi}\int_0^\pi\bigl(\sin(1+n)t +
\sin(1-n)t\bigr)\dd t
= \frac2\pi\cdot\frac{1 + \cos n\pi}{1 - n^2}
\]
pour $n \neq 1$ (et $a_1 = 0$ directement)~: nul pour $n$ impair, et
$a_{2k} = \frac{-4}{\pi(4k^2-1)}$. D'après le
\cref{thm:b2:fourier:parseval}~(1) la convergence est normale, et
\[
\abs{\sin t} = \frac{2}{\pi} - \frac{4}{\pi}
\sum_{k\geq1}\frac{\cos(2kt)}{4k^2 - 1}
\qquad (t \in \R) .
\]
Vérification intégrée en $t = 0$~: l'identité exige
$\sum_{k\geq1}\frac{1}{4k^2-1} = \frac12$, ce que le télescopage
confirme~:
\[
\sum_{k\geq1}\frac{1}{4k^2-1}
= \frac12\sum_{k\geq1}\Bigl(\frac{1}{2k-1} -
\frac{1}{2k+1}\Bigr) = \frac12 . \checkmark
\]
Éclairage final~: le \omterm{def:b2:reduction:eigen}{spectre} de $\abs{\sin}$ ne vit que sur les
fréquences \emph{paires} --- redresser un sinus double son contenu
fréquentiel, ce qui explique pourquoi les redresseurs pleine onde
bourdonnent à $100$ ou $120$ hertz, le double de la fréquence du
secteur.
\end{example}

\begin{example}[Parseval comme machine à calculer]\label{ex:b2:fourier:zetafourodd}
Parseval transforme les développements en séries numériques en
gros. Appliquons-le à $f(t) = \abs t$ (\cref{exo:b2:fourier:2}~: $a_0
= \pi$, $a_n = \frac{-4}{\pi n^2}$ pour $n$ impair, le reste nul)~:
\[
\frac{1}{2\pi}\int_{-\pi}^{\pi}t^2\,\dd t = \frac{\pi^2}{3}
= \frac{a_0^2}{4} + \frac12\sum_{n \text{ impair}} a_n^2
= \frac{\pi^2}{4} +
\frac{8}{\pi^2}\sum_{n\text{ impair}}\frac{1}{n^4} ,
\]
d'où
\[
\sum_{n\text{ impair}}\frac{1}{n^4}
= \frac{\pi^2}{8}\Bigl(\frac{\pi^2}{3} -
\frac{\pi^2}{4}\Bigr) = \frac{\pi^4}{96} .
\]
Recoupement avec l'\cref{exo:b2:fourier:3}~: en scindant
$\sum\frac1{n^4}$ en parties impaire et paire on obtient la
comptabilité à la $\zeta$ $S = S_{\mathrm{impair}} + \frac{S}{16}$,
donc $S = \frac{16}{15}\cdot\frac{\pi^4}{96} = \frac{\pi^4}{90}$ ---
exactement la valeur trouvée là-bas par une autre fonction. Deux
développements, un nombre~: la cohérence, c'est l'isométrie de
Parseval à l'œuvre. Éclairage final~: chaque nouveau développement
de Fourier est une machine à engendrer des identités de séries~; la
Partie~II du devoir explique pourquoi la machine ne peut jamais se
contredire.
\end{example}

\begin{omfigure}
\begin{tikzpicture}
\begin{axis}[omaxis, width=9.2cm, height=5.6cm,
    xmin=-3.6, xmax=3.6, ymin=-1.5, ymax=1.5,
    xtick={-3.1416, 0, 3.1416},
    xticklabels={$-\pi$, $0$, $\pi$}, ytick={-1,1}]
  \addplot[omExo, thick] coordinates {(-3.1416,-1) (0,-1)};
  \addplot[omExo, thick] coordinates {(0,1) (3.1416,1)};
  \addplot[omDef, thick, domain=-3.3:3.3, samples=200]
    {4/pi*sin(deg(x))};
  \addplot[omThm, very thick, domain=-3.3:3.3, samples=400]
    {4/pi*(sin(deg(x)) + sin(deg(3*x))/3 + sin(deg(5*x))/5 +
     sin(deg(7*x))/7 + sin(deg(9*x))/9)};
\end{axis}
\end{tikzpicture}

{\small Le signal carré (gris) et les \omterm{def:b2:fourier:coefficients}{sommes partielles de Fourier}
$S_1$ (bleu) et $S_9$ (rouge)~: convergence en tout point de
continuité, mais un dépassement persistant de $\sim 9\%$ près des
sauts --- le \emph{phénomène de Gibbs}\index{phénomène de Gibbs}. La
\omterm{def:b2:funcseq:def}{convergence uniforme} échoue précisément parce que la limite est
discontinue.}
\end{omfigure}

\begin{example}[Translation et modulation]\label{ex:b2:fourier:shiftrules}
Deux règles d'une ligne engendrent de nombreux développements à
partir d'un seul. Pour $a \in \R$, en substituant $s = t - a$~:
\[
c_n\bigl(f(\cdot - a)\bigr)
= \frac{1}{2\pi}\int_{-\pi}^{\pi}f(t - a)\eu^{-\iu nt}\dd t
= \eu^{-\iu na}\,c_n(f)
\qquad\text{(la translation module le spectre)},
\]
et directement à partir de la définition,
\[
c_n\bigl(\eu^{\iu kt}f\bigr) = c_{n-k}(f)
\qquad\text{(la modulation translate le spectre)}.
\]
Instance travaillée~: la dent de scie décalée de $\pi$, $g(t) = f(t
- \pi)$ avec $f(t) = t$, a pour coefficients $b_n$
$(-1)^n\cdot\frac{2(-1)^{n+1}}{n} = -\frac2n$~: le développement $g
\sim -2\sum\frac{\sin nt}{n}$ de la dent de scie qui saute en $0$ au
lieu de $\pi$ --- aucune intégrale recalculée. Éclairage final~: les
décalages temporels ne font que tourner les phases, jamais les
amplitudes ($\abs{c_n}$ est invariant par décalage), ce qui explique
pourquoi l'énergie (Parseval) et la classe de convergence sont des
propriétés de la \emph{forme} du signal, et non de l'endroit où
l'horloge démarre.
\end{example}

\begin{remark}[Pièges classiques]
\emph{(i) Trois convergences, trois monnaies~:} simple (Dirichlet~:
demande $C^1$ par morceaux, paie le \emph{milieu} à chaque saut ---
jamais la valeur unilatérale), uniforme (demande une limite \omterm{def:b2:metric:continuity}{continue}~;
impossible à travers un saut, Gibbs en est le symptôme visible), et
en moyenne quadratique (Parseval~: la plus robuste, aveugle aux points
individuels). Toujours nommer celle qu'on revendique. \emph{(ii) Pas
de dérivation terme à terme par défaut~:} dériver terme à terme la
série de la dent de scie de l'\cref{ex:b2:fourier:basel} donne $\sum
2(-1)^{n+1}\cos nt$, dont les termes ne tendent même pas vers $0$ ---
les théorèmes de transfert du chapitre sur les suites de fonctions
demandent la \omterm{def:b2:funcseq:def}{convergence uniforme} de la série \emph{dérivée}, que le
saut détruit. La régularité d'abord, la dérivation ensuite
(\cref{exo:b2:fourier:6} est le dictionnaire). \emph{(iii) Sommes
partielles \omterm{def:b2:quadratic:adjoint}{symétriques}~:} le théorème de Dirichlet concerne $S_N =
\sum_{-N}^{N}$~; réarranger ou sommer un côté d'abord peut changer
la divergence en convergence et inversement. \emph{(iv) Dérive de
normalisation~:} les conventions diffèrent d'un ouvrage à l'autre
($\frac{1}{2\pi}$ ou $\frac1\pi$ devant, période $2\pi$ ou $1$)~; les
invariants fiables sont les relations d'orthonormalité --- recalculer
$\langle e_m, e_n\rangle$ dans la convention en présence avant de se
fier à une formule.
\end{remark}

\begin{remark}[Où cela sert]
Parseval est le germe de la théorie $L^2$ des séries de Fourier~: le
volume de Licence~3 complète le tableau (les exponentielles forment
une base hilbertienne de $L^2$, et l'application $f \mapsto (c_n)$ est
une isométrie bijective). Au sein de ce volume, le devoir démontre le
\emph{théorème de Fejér} --- les moyennes de Cesàro de la série de
Fourier convergent \omterm{def:b2:funcseq:def}{uniformément} pour toute fonction $f$ \omterm{def:b2:metric:continuity}{continue}
périodique --- qui promeut le Parseval général admis en un théorème,
fournit le théorème de Weierstrass trigonométrique, et rapporte deux
dividendes spectaculaires~: le théorème d'équirépartition de Weyl et
l'inégalité isopérimétrique. Les mathématiques appliquées lisent ce
chapitre au quotidien~: \omterm{def:b2:reduction:eigen}{spectres} de signaux, harmoniques des systèmes
vibrants, et la transformée de Fourier rapide (dont l'avatar fini
était l'\cref{exo:b2:hermitian:10}).
\end{remark}

\begin{remark}[Perspectives dans ce volume]
Trois chapitres dialoguent avec celui-ci. En amont~: le chapitre
\omterm{def:b2:hermitian:adjoint}{hermitien} a fourni la géométrie (familles orthonormées, projections,
Bessel), et le chapitre sur les suites de fonctions l'analyse
(\omterm{def:b2:funcseq:def}{convergence uniforme}, théorèmes de transfert, identités approchées
--- le noyau de Fejér est aux séries de Fourier ce que les \omterm{thm:b2:funcseq:weierstrass}{polynômes
de Bernstein} étaient à Weierstrass). Latéralement~: la théorie au bord
du chapitre sur les séries entières revient via la sommation d'Abel,
réalisée ici par le noyau de Poisson (\cref{exo:b2:fourier:12}) --- le
rayon $r$ du disque jouant le rôle du paramètre de sommation. En
aval~: le chapitre sur les équations différentielles décompose un
forçage périodique en harmoniques et alimente chacune à la réponse en
fréquence de l'oscillateur~; la résonance survient lorsqu'un mode de
Fourier de l'entrée coïncide avec une fréquence propre, ce qui
explique pourquoi son devoir et celui de ce chapitre sont deux
moitiés d'une même histoire.
\end{remark}

\begin{omfigure}
\begin{tikzpicture}
\begin{axis}[omaxis, width=9.2cm, height=5.6cm,
    xmin=-3.6, xmax=3.6, ymin=-2.5, ymax=9.5,
    xtick={-3.1416, 0, 3.1416},
    xticklabels={$-\pi$, $0$, $\pi$}, ytick={8}]
  \addplot[omDef, thick, domain=0.02:3.3, samples=300]
    {sin(deg(8.5*x))/sin(deg(0.5*x))};
  \addplot[omDef, thick, domain=-3.3:-0.02, samples=300]
    {sin(deg(8.5*x))/sin(deg(0.5*x))};
  \addplot[omThm, very thick, domain=0.02:3.3, samples=300]
    {sin(deg(4*x))^2/(8*sin(deg(0.5*x))^2)};
  \addplot[omThm, very thick, domain=-3.3:-0.02, samples=300]
    {sin(deg(4*x))^2/(8*sin(deg(0.5*x))^2)};
  \node[font=\small, omDef] at (axis cs:1.05,5.2) {$D_8$};
  \node[font=\small, omThm] at (axis cs:-0.75,6.2) {$F_8$};
\end{axis}
\end{tikzpicture}

{\small Le \omterm{lem:b2:fourier:kernel}{noyau de Dirichlet} $D_8$ (bleu) oscille et prend des
valeurs négatives~; le noyau de Fejér $F_8$ (rouge) est positif, se
concentre en $0$, et a pour moyenne $1$~: une \emph{identité
approchée}. La positivité est exactement ce qui manque au \omterm{lem:b2:fourier:kernel}{noyau de
Dirichlet}, et ce qui rend le théorème de Fejér du devoir
inconditionnel.}
\end{omfigure}

\section{Exercices}

\begin{exercise}[$\star$]\label{exo:b2:fourier:1}
Calculer les \omterm{def:b2:fourier:coefficients}{coefficients de Fourier} du signal carré ($f = -1$ sur
$\intoo{-\pi}{0}$, $+1$ sur $\intoo{0}{\pi}$), énoncer la conclusion
de Dirichlet en $t = \frac\pi2$ et au saut $t = 0$, et retrouver la
série de Leibniz.
\end{exercise}

\begin{exercise}[$\star$]\label{exo:b2:fourier:2}
Développer $f(t) = \abs t$ ($\abs t \leq \pi$, $2\pi$-périodique) en
série de Fourier~; justifier la convergence \emph{normale}~; évaluer
en $t = 0$ pour obtenir $\sum_{k\geq0}\frac{1}{(2k+1)^2} =
\frac{\pi^2}{8}$, et en redéduire Bâle.
\end{exercise}

\begin{exercise}[$\star$]\label{exo:b2:fourier:3}
Développer $f(t) = t^2$ ($\abs t \leq \pi$) et en déduire
\[
\sum_{n\geq1}\frac{(-1)^{n+1}}{n^2} = \frac{\pi^2}{12},
\qquad
\sum_{n\geq1}\frac{1}{n^4} = \frac{\pi^4}{90}
\quad\text{(Parseval)}.
\]
\end{exercise}

\begin{exercise}[$\star\star$]\label{exo:b2:fourier:4}
Soit $f$ \omterm{def:b2:metric:continuity}{continue} $2\pi$-périodique avec $c_n(f) = 0$ pour tout $n$.
Prouver $f = 0$ \emph{(Parseval --- pour quelle classe est-il démontré
ici~? justifier que continuité plus $C^1$ par morceaux peut être
abandonné en admettant le Parseval général, ou donner l'argument de
densité en esquisse)}.
\end{exercise}

\begin{exercise}[$\star\star$]\label{exo:b2:fourier:5}
Pour $\alpha \notin \Z$, développer $f(t) = \cos(\alpha t)$ ($\abs t
\leq \pi$) et en déduire le développement en éléments simples de la
cotangente~:
\[
\pi\cot(\pi\alpha) = \frac{1}{\alpha} + \sum_{n\geq1}
\frac{2\alpha}{\alpha^2 - n^2} .
\]
\end{exercise}

\begin{exercise}[$\star\star$]\label{exo:b2:fourier:6}
Prouver que si $f$ est $2\pi$-périodique et $C^k$ avec $f^{(k)}$
\omterm{def:b2:metric:continuity}{continue} par morceaux, alors $c_n(f) = O\bigl(\abs n^{-k}\bigr)$~:
régularité du signal $=$ décroissance de son \omterm{def:b2:reduction:eigen}{spectre}.
\end{exercise}

\begin{exercise}[$\star\star\star$]\label{exo:b2:fourier:7}
(Inégalité de Wirtinger) Soit $f$ de classe $C^1$, $2\pi$-périodique,
avec $\int_{-\pi}^{\pi} f = 0$. Prouver
\[
\int_{-\pi}^{\pi} \abs{f}^2 \leq \int_{-\pi}^{\pi} \abs{f'}^2 ,
\]
avec égalité si et seulement si $f(t) = a\cos t + b\sin t$.
\emph{(Parseval des deux côtés~; comparer $\abs{c_n}^2$ et
$n^2\abs{c_n}^2$.)}
\end{exercise}

\begin{exercise}[$\star\star\star$]\label{exo:b2:fourier:8}
(La constante de Gibbs) Pour le signal carré de
l'\cref{exo:b2:fourier:1}, évaluer la somme partielle en $x_N =
\frac{\pi}{2N}$~: en posant $u_k = \frac{(2k+1)\pi}{2N}$ et $\Delta u
= \frac{\pi}{N}$, montrer que
\[
S_{2N-1}\Bigl(\frac{\pi}{2N}\Bigr)
= \frac{4}{\pi}\sum_{k=0}^{N-1} \frac{\sin u_k}{2k+1}
= \frac{2}{\pi}\sum_{k=0}^{N-1} \frac{\sin u_k}{u_k}\,\Delta u
\xrightarrow[N\to\infty]{}
\frac{2}{\pi}\int_0^{\pi}\frac{\sin u}{u}\,\dd u \approx 1.179 :
\]
une somme de Riemann de $\frac{2}{\pi}\cdot\frac{\sin u}{u}$ sur
$\intcc{0}{\pi}$ aux milieux. Conclure que le dépassement au-delà de
la valeur de saut $1$ ne s'annule pas quand $N \to \infty$.
\end{exercise}

\begin{exercise}[$\star$]\label{exo:b2:fourier:9}
Développer $\cos^3 t$ et $\sin^2 t\,\cos t$ en série de Fourier
\emph{(linéariser~; un polynôme trigonométrique est sa propre série
de Fourier, par unicité des coefficients)}. Quels sont $c_n$, $a_n$,
$b_n$ pour chacun~?
\end{exercise}

\begin{exercise}[$\star\star$]\label{exo:b2:fourier:10}
Soit $a > 0$ et $f(t) = \eu^{at}$ sur $\intoc{-\pi}{\pi}$, prolongée
de manière $2\pi$-périodique. Calculer
\[
c_n(f) = \frac{(-1)^n\sinh(a\pi)}{\pi(a - \iu n)} ,
\]
appliquer le théorème de Dirichlet au saut $t = \pi$, et en déduire
le développement en éléments simples de la cotangente hyperbolique~:
\[
\coth(\pi a) = \frac{1}{\pi a} + \sum_{n\geq1}
\frac{2a}{\pi(a^2 + n^2)} .
\]
\end{exercise}

\begin{exercise}[$\star\star$]\label{exo:b2:fourier:11}
(Convolution) Pour $f, g$ \omterm{def:b2:metric:continuity}{continues} et $2\pi$-périodiques on définit
\[
(f * g)(x) = \frac{1}{2\pi}\int_{-\pi}^{\pi} f(x - t)\,g(t)\,
\dd t .
\]
Montrer que $f * g = g * f$, que $c_n(f*g) = c_n(f)\,c_n(g)$
\emph{(pour échanger les deux intégrales d'un intégrande \omterm{def:b2:metric:continuity}{continu},
comparer les deux fonctions de la borne supérieure~: toutes deux
s'annulent au bord gauche et ont la même dérivée, par continuité et
dérivation sous le signe intégral)}, et que $S_N(f) = f * D_N$ pour le
\omterm{lem:b2:fourier:kernel}{noyau de Dirichlet}. (Les moyennes de Fejér du devoir sont de même des
convolutions, $\sigma_N(f) = f * F_N$.)
\end{exercise}

\begin{exercise}[$\star\star\star$]\label{exo:b2:fourier:12}
(Noyau de Poisson~: moyennes d'Abel des séries de Fourier) Pour $0
\leq r < 1$ on pose $P_r(t) = \sum_{n\in\Z} r^{\abs n}\eu^{\iu nt}$.
\begin{enumerate}
  \item Sommer les deux séries géométriques et montrer
        \[
        P_r(t) = \frac{1 - r^2}{1 - 2r\cos t + r^2} > 0,
        \qquad
        \frac{1}{2\pi}\int_{-\pi}^{\pi}P_r = 1 .
        \]
  \item Montrer que pour $\delta \leq \abs t \leq \pi$~: $P_r(t)
        \leq \frac{1 - r^2}{1 - 2r\cos\delta + r^2} \to 0$ quand
        $r \to 1^-$, \omterm{def:b2:funcseq:def}{uniformément}.
  \item En déduire que pour toute fonction $f$ \omterm{def:b2:metric:continuity}{continue}
        $2\pi$-périodique, les \emph{moyennes d'Abel} $(f *
        P_r)(x) = \sum_n r^{\abs n}c_n(f)\,\eu^{\iu nx}$ convergent
        vers $f$ \omterm{def:b2:funcseq:def}{uniformément} quand $r \to 1^-$ --- la sœur à
        paramètre \omterm{def:b2:metric:continuity}{continu} du théorème de Fejér, et l'incarnation de
        Fourier de la sommation d'Abel du chapitre sur les séries
        entières.
\end{enumerate}
\end{exercise}

\section{Problème~: le théorème de Fejér et ses dividendes}

\begin{problem}\label{pb:b2:fourier:1}
Le théorème de Dirichlet demande $f$ $C^1$ par morceaux~; pour $f$
seulement \omterm{def:b2:metric:continuity}{continue} les sommes partielles $S_N(f)$ peuvent mal se
comporter. La découverte de Fejér~: leurs \emph{moyennes de Cesàro}
ne le font jamais. Le moteur est la positivité du noyau de Fejér, et
la moisson est immense~: approximation trigonométrique uniforme
(Weierstrass), unicité des \omterm{def:b2:fourier:coefficients}{coefficients de Fourier}, Parseval pour
toute fonction de ce chapitre (levant l'« admis » du
\cref{thm:b2:fourier:parseval}), le théorème d'équirépartition de
Weyl, et --- couronnant un siècle de géométrie --- l'inégalité
isopérimétrique.\index{théorème de Fejér}\index{équirépartition}\index{inégalité isopérimétrique}
Dans tout le problème, $f$ est $2\pi$-périodique et \omterm{def:b2:metric:continuity}{continue} par
morceaux, et
\[
\sigma_N(f) = \frac{S_0(f) + S_1(f) + \dots +
S_{N-1}(f)}{N} .
\]

\medskip
\noindent\textbf{Partie I --- Le noyau de Fejér.}
\begin{enumerate}
  \item Montrer que $\sigma_N(f)(x) =
        \frac{1}{2\pi}\int_{-\pi}^{\pi}f(x+u)\,F_N(u)\,\dd u$
        avec $F_N = \frac{D_0 + \dots + D_{N-1}}{N}$, et que
        $\frac{1}{2\pi}\int_{-\pi}^{\pi}F_N = 1$.
  \item Prouver la forme fermée, pour $u \notin 2\pi\Z$~:
        \[
        F_N(u) = \frac{1}{N}\,
        \frac{\sin^2\bigl(\frac{Nu}{2}\bigr)}
        {\sin^2\bigl(\frac u2\bigr)} \;\geq\; 0
        \]
        \emph{(sommer $\sin\bigl((n+\frac12)u\bigr)$ comme la
        partie imaginaire d'une série géométrique)}.
  \item Montrer l'estimation de concentration~: pour $0 < \delta
        \leq \abs u \leq \pi$,
        \[
        F_N(u) \leq \frac{1}{N\sin^2\frac\delta2}
        \xrightarrow[N\to\infty]{} 0
        \quad\text{uniformément} :
        \]
        $(F_N)$ est une identité approchée positive.
  \item (Théorème de Fejér) Prouver~: si $f$ est \omterm{def:b2:metric:continuity}{continue} et
        $2\pi$-périodique, alors $\sigma_N(f) \to f$
        \emph{\omterm{def:b2:funcseq:def}{uniformément}} sur $\R$ \emph{(scinder l'intégrale de
        $\bigl(f(x+u) - f(x)\bigr)F_N(u)$ en $\abs u =
        \delta$~; utiliser Heine et les questions 1--3)}.
  \item Pour $f$ \omterm{def:b2:metric:continuity}{continue} par morceaux, montrer la version
        ponctuelle $\sigma_N(f)(x) \to \frac{f(x^+) +
        f(x^-)}{2}$ en tout $x$, et la majoration uniforme
        $\norm{\sigma_N(f)}_\infty \leq \norm f_\infty$
        \emph{(positivité~!)}.
\end{enumerate}

\medskip
\noindent\textbf{Partie II --- Weierstrass, unicité, Parseval.}
\begin{enumerate}[resume]
  \item (Weierstrass trigonométrique) En déduire~: toute fonction
        \omterm{def:b2:metric:continuity}{continue} $2\pi$-périodique est limite uniforme de
        polynômes trigonométriques.
  \item (Unicité) En déduire~: une fonction $f$ \omterm{def:b2:metric:continuity}{continue} avec
        $c_n(f) = 0$ pour tout $n$ est identiquement nulle --- deux
        fonctions \omterm{def:b2:metric:continuity}{continues} périodiques ayant les mêmes
        \omterm{def:b2:fourier:coefficients}{coefficients de Fourier} coïncident (\cref{exo:b2:fourier:4},
        maintenant sans admission).
  \item (Parseval, cas \omterm{def:b2:metric:continuity}{continu}) En utilisant la propriété de
        projection de $S_N$ (\cref{prop:b2:fourier:bessel}) et
        $\sigma_Nf \in \mathcal T_{N-1} \subseteq \mathcal
        T_N$, prouver
        \[
        \norm{f - S_Nf}_2 \leq \norm{f - \sigma_Nf}_2
        \leq \norm{f - \sigma_Nf}_\infty
        \xrightarrow[N\to\infty]{} 0 ,
        \]
        et conclure l'identité de Parseval pour toute fonction $f$
        \emph{\omterm{def:b2:metric:continuity}{continue}} $2\pi$-périodique.
  \item (Parseval, cas \omterm{def:b2:metric:continuity}{continu} par morceaux) Étant donnée $f$
        \omterm{def:b2:metric:continuity}{continue} par morceaux et $\varepsilon > 0$, construire une
        fonction $g$ \omterm{def:b2:metric:continuity}{continue} périodique avec $\norm{f - g}_2 \leq
        \varepsilon$ \emph{(remplacer $f$ par une interpolation
        affine sur de minuscules intervalles autour des sauts)}, et
        en déduire $\norm{f - S_Nf}_2 \to 0$
        \emph{(utiliser Bessel~: $\norm{S_Nh}_2 \leq \norm h_2$)}~:
        Parseval vaut dans toute la généralité énoncée au
        \cref{thm:b2:fourier:parseval} --- l'« admis » a disparu.
  \item (Pas de Gibbs pour Fejér) Contraster avec
        l'\cref{exo:b2:fourier:8}~: montrer que pour le signal carré
        $f$, $\abs{\sigma_N(f)} \leq 1$ partout, pour tout $N$ ---
        la moyennisation de Cesàro efface le dépassement qui hante
        $S_N$. Expliquer en une phrase quelle propriété de $F_N$ en
        est responsable.
\end{enumerate}

\medskip
\noindent\textbf{Partie III --- Vitesses.}
\begin{enumerate}[resume]
  \item Prouver les deux majorations du noyau, pour $0 < \abs u
        \leq \pi$~:
        \[
        F_N(u) \leq N,
        \qquad
        F_N(u) \leq \frac{\pi^2}{N u^2}
        \]
        \emph{(pour la première, $\abs{\sin N\theta} \leq
        N\abs{\sin\theta}$ par récurrence~; pour la seconde,
        $\sin\frac u2 \geq \frac{u}{\pi}$ sur
        $\intcc{0}{\pi}$)}.
  \item En déduire l'estimation du premier moment
        \[
        \frac{1}{2\pi}\int_{-\pi}^{\pi}\abs u\,F_N(u)\,\dd u
        \;\leq\; \frac{C\,\ln N}{N}
        \qquad (N \geq 2)
        \]
        pour une constante explicite \emph{(scinder en $\abs u =
        \frac1N$)}.
  \item Conclure~: si $f$ est $L$-lipschitzienne et
        $2\pi$-périodique, alors
        \[
        \norm{\sigma_N f - f}_\infty \leq \frac{C\,L\ln N}{N} .
        \]
  \item (Saturation) Calculer $\sigma_N(e_1)$ pour $e_1(t) =
        \eu^{\iu t}$ et montrer $\norm{\sigma_Ne_1 - e_1}_\infty
        = \frac1N$~: même pour les fonctions les plus régulières,
        Fejér ne \omterm{def:b2:integration:improper}{converge} pas plus vite que $\frac1N$ --- l'exact
        analogue de la saturation de Bernstein dans le devoir du
        chapitre sur les suites de fonctions.
  \item (Localisation) Montrer~: si $f$ (\omterm{def:b2:metric:continuity}{continue} par morceaux)
        s'annule sur $\intoo{x - \delta}{x + \delta}$, alors
        $\sigma_N(f)(x) \to 0$, aussi sauvage que soit $f$
        ailleurs --- la convergence des moyennes en $x$ ne voit que
        $f$ près de $x$.
\end{enumerate}

\medskip
\noindent\textbf{Partie IV --- Le théorème d'équirépartition de Weyl.}
Une suite $(x_n)_{n\geq1}$ de $\intco{0}{1}$ est
\emph{équirépartie} lorsque, pour tout intervalle
$\intcc{a}{b} \subseteq \intcc{0}{1}$,
\[
\frac{\#\{n \leq N : x_n \in \intcc ab\}}{N}
\xrightarrow[N\to\infty]{} b - a .
\]
\begin{enumerate}[resume]
  \item Montrer que $(x_n)$ est équirépartie dès que
        $\frac1N\sum_{n\leq N}f(x_n) \to \int_0^1 f$ pour toute
        fonction $f$ \emph{\omterm{def:b2:metric:continuity}{continue}} $1$-périodique \emph{(encadrer
        l'indicatrice de $\intcc ab$ entre deux fonctions
        \omterm{def:b2:metric:continuity}{continues} affines par morceaux dont les intégrales diffèrent
        de $\varepsilon$)}.
  \item (Critère de Weyl, suffisance) Supposons
        \[
        \frac{1}{N}\sum_{n=1}^{N}\eu^{2\iu\pi kx_n}
        \xrightarrow[N\to\infty]{} 0
        \qquad\text{pour tout } k \in \Z\setminus\{0\} .
        \]
        Montrer $\frac1N\sum f(x_n) \to \int_0^1f$ d'abord pour les
        polynômes trigonométriques, puis pour toute fonction $f$
        \omterm{def:b2:metric:continuity}{continue} $1$-périodique par la question 6 (transportée à la
        période $1$)~: avec la question 16, $(x_n)$ est
        équirépartie.
  \item Soit $\alpha$ irrationnel et $x_n = \{n\alpha\}$ (partie
        fractionnaire). Majorer la somme géométrique
        \[
        \Bigl|\sum_{n=1}^{N}\eu^{2\iu\pi kn\alpha}\Bigr|
        \leq \frac{2}{\abs{1 - \eu^{2\iu\pi k\alpha}}}
        \qquad (k \neq 0),
        \]
        et conclure le \emph{théorème de Weyl}~: $(\{n\alpha\})$
        est équirépartie dans $\intco{0}{1}$.
  \item En déduire que $(\{n\alpha\})$ est dense dans
        $\intcc{0}{1}$ pour $\alpha$ irrationnel, et expliquer en
        une phrase pourquoi l'équirépartition est strictement plus
        forte que la densité.
  \item (Premiers chiffres) Prouver que la proportion d'entiers
        $n \leq N$ tels que $2^n$ ait pour premier chiffre
        (décimal) $1$ tend vers $\log_{10}2 \approx 0.301$
        \emph{(premier chiffre $1$ signifie $\{n\log_{10}2\} \in
        \intco{0}{\log_{10}2}$~; montrer que $\log_{10}2$ est
        irrationnel)}.
\end{enumerate}

\medskip
\noindent\textbf{Partie V --- L'inégalité isopérimétrique.} Soit
$\Gamma$ une courbe fermée simple $C^1$ de longueur $L$ enfermant
une aire algébrique $A$, paramétrée par longueur d'arc mise à
l'échelle~: $z(t) = x(t) + \iu y(t)$, $2\pi$-périodique, avec
$\abs{z'(t)} = \frac{L}{2\pi}$ constant~; l'aire enfermée est
\[
A = \frac12\int_0^{2\pi}\bigl(x\,y' - y\,x'\bigr)\dd t
= \frac{1}{2}\,\Im\int_0^{2\pi}\conj{z}\,z'\,\dd t
\]
(prise ici comme définition de l'aire algébrique~; le chapitre sur
les intégrales multiples démontre qu'elle coïncide avec l'aire
intuitive, via la formule de Green).
\begin{enumerate}[resume]
  \item Développer $z(t) = \sum_{n\in\Z}c_n\eu^{\iu nt}$ (la série
        d'une fonction $C^1$, \omterm{def:b2:funcseq:series}{normalement} convergente) et prouver,
        par Parseval appliqué à $z'$~:
        \[
        \frac{L^2}{2\pi} = \int_0^{2\pi}\abs{z'}^2\dd t
        = 2\pi\sum_{n\in\Z}n^2\abs{c_n}^2 .
        \]
  \item Prouver de même $A = \pi\sum_{n\in\Z}
        n\,\abs{c_n}^2$ \emph{(Parseval sous sa forme polarisée~:
        $\frac{1}{2\pi}\int\conj f g = \sum
        \conj{c_n(f)}c_n(g)$, appliqué à $f = z$, $g = z'$)}.
  \item (Hurwitz) Conclure~:
        \[
        L^2 - 4\pi A = 4\pi^2\sum_{n\in\Z}
        (n^2 - n)\abs{c_n}^2 \;\geq\; 0 ,
        \]
        avec égalité si et seulement si $z(t) = c_0 +
        c_1\eu^{\iu t}$ --- un cercle. L'inégalité isopérimétrique~:
        parmi les courbes fermées de longueur $L$, seul le cercle
        enferme l'aire $\frac{L^2}{4\pi}$.
  \item Vérifications de cohérence~: vérifier l'égalité pour le
        cercle de rayon $R$ et l'inégalité stricte pour le carré de
        côté $a$~; expliquer pourquoi $n^2 - n \geq 0$ pour tout
        entier $n$, y compris les négatifs, et où la vitesse
        constante du paramétrage a servi.
  \item Synthèse. En une phrase chacune~: (i) l'unique propriété de
        $F_N$ d'où découlent les Parties I--III, et qui manque à
        $D_N$~; (ii) comment la sommation de Cesàro relie ici au
        devoir du chapitre sur les séries entières (Frobenius)~;
        (iii) quel dividende n'a utilisé que Weierstrass (question
        6) et lequel a demandé Parseval \omterm{def:b2:metric:complete}{complet}~; (iv) une phrase
        sur ce qu'ajoute le volume de Licence~3 (complétude $L^2$~:
        les séries de Fourier comme base hilbertienne).
\end{enumerate}
\end{problem}
